\documentclass[11pt,a4paper]{article}
\usepackage[margin=0.85in,top=1in,bottom=1in]{geometry}
\usepackage{amsmath,amssymb,amsfonts}
\usepackage{graphicx}
\usepackage{float}
\usepackage{enumitem}
\usepackage{microtype}
\usepackage{caption}
\usepackage[hidelinks]{hyperref}
\usepackage[most,breakable]{tcolorbox}
\tcbuselibrary{breakable,skins}
\usepackage{fontspec}
\setmainfont{DejaVu Serif}
\setsansfont{DejaVu Sans}
\setmonofont{DejaVu Sans Mono}
\captionsetup{font=small,labelfont=bf,skip=4pt}
\setlist{leftmargin=1.5em,topsep=2pt}

%% Breakable problem card — prevents content cutoff across pages
\newtcolorbox{solcard}[3]{
  breakable, enhanced,
  colback=white, colframe=gray!50, arc=2pt,
  title={\textbf{#1}\hfill\textsf{\small #3\,pts}},
  fonttitle=\normalsize\sffamily\bfseries,
  before upper={\small\textit{#2}\par\smallskip\hrule height 0.4pt\smallskip},
  top=4pt, bottom=6pt, left=7pt, right=7pt,
  parbox=false,
}

\title{\textbf{Measuring the radius of a star --- Solution}}
\author{\normalsize Y23 (2023) — English translation}
\date{}
\begin{document}\maketitle\vspace{-0.5em}
\begin{solcard}{A1}{Let us consider the case when rays 1 and 2 are parallel to the $z$ axis, and mirror M3 is parallel to the $xy$ plane. Find the difference in optical paths $\Delta_{12}$ between rays 1 and 2 if they both fall on the screen at point $O$.}{0.20}

In horizontal sections, the path difference is $L_1 - L_2$, and the first ray also travels a distance $L_3$ twice. Thus, the total difference is: $$\Delta_{12} = L_1+2L_3-L_2$$

\par\smallskip\noindent\textbf{Answer:} \textbackslash{}textbf{Answer:} \[\Delta_{12}=L_1+2L_3-L_2\]
\end{solcard}

\begin{solcard}{A2}{Let us consider the case when the mirror M3 is deflected by an angle $\alpha$ from the plane $xy$ by rotating around the axis $x$ (see Fig. 1). If ray 1' hits the screen at point $O$, find the optical path difference $\Delta_{1'2}$ between rays 1' and 2. Show that, to a first approximation, the path difference is independent of $\alpha$ and simplify your answer at $\alpha\ll1$.}{1.40}

Note that the triangle obtained from the paths of rays 1 and 1' from the mirror M3 to the point $O$ is isosceles. Also, changes in the horizontal and vertical paths due to changes in the point of incidence of the beam compensate each other. The only parameter affecting the path traveled by beam 1 is the coordinate $z$ of the point of impact on the mirror M3: it affects the distance traveled from HM to M3. But the optical path of beam 2 does not change in any way. We get: $$\Delta_{1'2} = \Delta_{12} - \left(2(L_3+L_4)\sin{\alpha}\right)\cdot \sin{\alpha} = \Delta_{12} - 2(L_3+L_4)\sin^2{\alpha}$$Substituting result прошлого partа: $$\Delta_{12'} = L_1+2L_3-L_2 - 2(L_3+L_4)\sin^2{\alpha} \approx L_1+2L_3-L_2 - 2(L_3+L_4)\alpha^2$$

\par\smallskip\noindent\textbf{Answer:} \textbackslash{}textbf{Answer:} \[\Delta_{1'2}=L_1+2L_3-L_2-(L_3+L_4)\frac{\sin\alpha\sin2\alpha}{\cos\alpha}=L_1+2L_3-L_2-2(L_3+L_4)\alpha^2\]
\end{solcard}

\begin{solcard}{B1}{Let rays 1 and 2 fall at fixed points of mirrors M1 and M2, respectively. Show in the solution sheets that, to a first approximation, the optical paths of rays 1 and 2 from the points of reflection from the mirrors M1 and M2, respectively, to the screen do not depend on the angles $\theta$ and $\alpha$.}{0.80}
\end{solcard}

\begin{solcard}{B2}{Let the angles $\alpha$ and $\theta$ be fixed, and the positions of the points of incidence of rays 1 and 2 on the mirrors M1 and M2, respectively, can change. Show in the solution sheets that the optical paths of rays 1 and 2 from the dotted line corresponding to the coordinate $z_\text{п}=L_3+L_4$ (meaning the points of intersection of the dotted line with the trajectories of rays 1 and 2 before falling on the mirrors M1 and M2, respectively) to the screen, to a first approximation, do not depend on the angles $\theta$ and $\alpha$.}{0.80}
\end{solcard}

\begin{solcard}{B3}{Express $b$ and $\delta_O$ in terms of $\lambda$, $L_1$, $L_2$, $L_3$, $L_4$.}{1.00}

$\delta_O$ can be found using the previous part of the problem, since this is the value of the phase difference for $\theta = 0$. Note that the phase difference arises not only due to the difference in optical paths, but also due to reflections: beam 1 is reflected one time more, so it is additionally shifted by $\pi$. $$\delta_O = \cfrac{2\pi}{\lambda}\left(L_1+2L_3-L_2\right) + \pi$$Можно заметить, что в силу утверждений, доказанных в B1 and B2 for/at изменении угла $\theta$ сность фаз возникает только from-за изменения путей до пересечения с partandрной линией $z_\text{п} = L_3 + L_4$, that is: $$\cfrac{2\pi}{\lambda}(L_1 + L_2) \theta = \cfrac{2\pi}{\lambda}b\theta$$$$b = L_1+L_2$$

\par\smallskip\noindent\textbf{Answer:} \textbackslash{}textbf{Answer:} \[b=L_1+L_2\\\delta_O=\frac{2\pi}\lambda\big[L_1+2L_3-L_2\big]+\pi\]
\end{solcard}

\begin{solcard}{B4}{Show in the solution sheets that $$\delta_y=\delta_O-\cfrac{2\pi}{\lambda}\cdot 2\alpha y $$ Note: Further, in all points of the problem, you can use the statements of points B1-B4, even if they are not proven.}{0.60}

Note: Further, in all points of the problem, you can use the statements of points B1-B4, even if they are not proven.


Consider two rays 1 that arrive at points on the screen located at a distance $y$. Due to the smallness of $\alpha$, the distance between the points of reflection of these rays from M3 is also equal to $y$. Note that the difference in optical paths between them arises only because of the different coordinates $z$ of the reflection points from M3: the path from HM to M3 and from M3 to the screen is shortened. $$\Delta L = -2y\alpha$$$$\delta_y = \delta_O - \cfrac{2\pi}{\lambda}\cdot 2\alpha y$$

\par\smallskip\noindent\textbf{Answer:} \textbackslash{}textbf{Answer:} \[\delta_y=\frac{2\pi}\lambda\big[L_1+2L_3-L_2-2\alpha y\big]+\pi\]
\end{solcard}

\begin{solcard}{B5}{Thus, an interference pattern will appear on the screen. Find the distance $y_1$ to the interference maximum closest to point $O$. Also find the distance $\Delta y_m$ between two adjacent interference maxima. Give numerical answers in the case of $L_1=0.50~\mathrm{m}$, $L_2=1.50~\mathrm{m}$, $L_3=0.50~\mathrm{m}$, $L_4=1.00~\mathrm{m}$, $\alpha=10^{-4}$, $\theta=10^{-5}$, $\lambda=500~\mathrm{nm}$.}{0.60}

Substituting numerical values, we get: $$\delta = 81\pi - 800\cfrac{\pi}{1\text{m}}\cdot y$$Максимум соответсвует $\delta = 2\pi n$. At the point $O$ находится минимум. $$y_1 = \cfrac{\lambda}{4\alpha} = 1.25\text{mm}$$$$\Delta y_m = \cfrac{\lambda}{2\alpha} = 2.5\text{mm}$$

\par\smallskip\noindent\textbf{Answer:} \textbackslash{}textbf{Answer:} \[y_1=1.25\cdot10^{-3}~\mathrm{m}\\\Delta y_m=\frac\lambda{2\alpha}=2.5~\mathrm{mm}\]
\end{solcard}

\begin{solcard}{B6}{Find how the intensity $I(y)$ on the screen depends on the coordinate $y$ if its maximum value is $I_0$, substituting here the numerical values ​​from the previous paragraph. Consider the effect of the beam splitter on the intensity of transmitted light.}{0.90}

The field amplitude in beam 1 is $\sqrt{2}$ times less than in the second. Then the intensity is proportional to the square of the amplitude: $$I \sim 1 + 2 + 2\sqrt{2}\cos{\delta}$$$$I(y)=I_0\frac{3-2\sqrt2\cos\left[80\pi\left(1-10\cfrac{y}{1\text{m}}\right)\right]}{3+2\sqrt2}$$

\par\smallskip\noindent\textbf{Answer:} \textbackslash{}textbf{Answer:} \[I(y)=I_0\frac{3-2\sqrt2\cos\left[80\pi\left(1-10\cfrac{y}{1\text{m}}\right)\right]}{3+2\sqrt2}\]
\end{solcard}

\begin{solcard}{C1}{Express the wave vector $\vec k$ of the incident rays in terms of $\theta_x$, $\theta_y$, the modulus of the wave vector $k=2\pi/\lambda$ and the unit unit vectors of the coordinate axes.}{0.40}

Due to the smallness of the angles $\theta_x$ and $\theta_y$: $$k_z = - k$$$$k_x = -k\theta_x$$$$k_y = -k\theta_y$$

\par\smallskip\noindent\textbf{Answer:} \textbackslash{}textbf{Answer:} \[\vec k=-k\big(\theta_x\hat x+\theta_y\hat y+\hat z\big)\]
\end{solcard}

\begin{solcard}{C2}{Express the phase difference between beams 1'' and 2'' in terms of $\lambda$, $L_1$, $L_2$, $L_3$, $L_4$, $\theta_x$, $\theta_y$, $\alpha$, $y$.}{0.60}

Note that the initial displacement along the $x$ axis of the rays hitting one point on the screen is of the order of $L_3 \theta_x$, the difference in the optical paths is correspondingly proportional to $\theta_x^2$, so the contribution $\theta_x$ can be ignored. The answer is similar to the previous part of the problem: \[\delta_{1''2''}=\cfrac{2\pi}{\lambda}\big((L_1+L_2)\theta_y+(L_1+2L_3-L_2)-2\alpha y\big)+\pi\]

\par\smallskip\noindent\textbf{Answer:} \textbackslash{}textbf{Answer:} \[\delta_{1''2''}=\cfrac{2\pi}{\lambda}\big((L_1+L_2)\theta_y+(L_1+2L_3-L_2)-2\alpha y\big)+\pi\]
\end{solcard}

\begin{solcard}{D1}{Write $V$ as an integral over one variable. It is not necessary to perform integration explicitly. Hint: the integral of some function $f(\theta_y)$, depending only on $\theta_y$, over the surface of a disk with radius $\theta_R$ can be represented as: \[I=\int\limits_{-\theta_R}^{\theta_R}\mathrm d\theta_y\int\limits_{-\sqrt{\theta^2_R-\theta^2_y}}^{\sqrt{\theta^2_R-\theta^2_y}}\mathrm d\theta_x~f(\theta_y)=\int\limits_{-\theta_R}^{\theta_R}\mathrm d\theta_y\cdot 2\sqrt{\theta^2_R-\theta^2_y}f(\theta_y).\]}{1.80}

Hint: the integral of some function $f(\theta_y)$, depending only on $\theta_y$, over the surface of a disk with radius $\theta_R$ can be represented as: \[I=\int\limits_{-\theta_R}^{\theta_R}\mathrm d\theta_y\int\limits_{-\sqrt{\theta^2_R-\theta^2_y}}^{\sqrt{\theta^2_R-\theta^2_y}}\mathrm d\theta_x~f(\theta_y)=\int\limits_{-\theta_R}^{\theta_R}\mathrm d\theta_y\cdot 2\sqrt{\theta^2_R-\theta^2_y}f(\theta_y).\]


Parts of the star emit independently, so the intensities from them can be added. $$dI = J d\theta_x d\theta_y\left(3 - 2\sqrt{2}\cos{\left(\frac{2\pi}\lambda\big(L_1+2L_3-L_2-2\alpha y\big)+\cfrac{2\pi}{\lambda}(L_1 + L_2 + d) \theta_y\right)}\right)$$Denote $\varphi_0 = \frac{2\pi}\lambda\big(L_1+2L_3-L_2\big)$. $$I = \int\limits_{-\theta_R}^{\theta_R} J\sqrt{\theta^2_R - \theta^2_y}\left(3 - 2\sqrt{2}\cos{\left(-\frac{4\pi}\lambda\alpha y+\cfrac{2\pi}{\lambda}(L_1 + L_2+d) \theta_y + \varphi_0\right)}\right)d\theta_y$$Note that: $$I_{max} = \int\limits_{-\theta_R}^{\theta_R} J\sqrt{\theta^2_R - \theta^2_y}\left(3 + 2\sqrt{2}\cos{\left(\cfrac{2\pi}{\lambda}(L_1 + L_2 + d) \theta_y\right)}\right)d\theta_y$$$$I_{min} = \int\limits_{-\theta_R}^{\theta_R} J\sqrt{\theta^2_R - \theta^2_y}\left(3 - 2\sqrt{2}\cos{\left(\cfrac{2\pi}{\lambda}(L_1 + L_2 + d) \theta_y\right)}\right)d\theta_y$$Then видность: $$V = \cfrac{2\sqrt{2}\int\limits_{-\theta_R}^{\theta_R} \sqrt{\theta^2_R - \theta^2_y}\cos{\left(\cfrac{2\pi}{\lambda}(L_1 + L_2 + d) \theta_y\right)}d\theta_y}{3\int\limits_{-\theta_R}^{\theta_R} \sqrt{\theta^2_R - \theta^2_y}d\theta_y}$$Denote $\omega = $Note that the integrable function is even, therefore: \[V=\frac{8\sqrt2}{3\pi}\int\limits_0^1\mathrm dw\sqrt{1-w^2}\cos(\eta w)\]

\par\smallskip\noindent\textbf{Answer:} \textbackslash{}textbf{Answer:} \[V=\frac{8\sqrt2}{3\pi}\int\limits_0^1\mathrm dw\sqrt{1-w^2}\cos(\eta w)\]
\end{solcard}

\begin{solcard}{D2}{Find Betelgeuse's angular radius $\theta_R$. Find its radius $R$ if the distance to it is $h=6.079\cdot10^{15}~\mathrm{km}$. Hint: The first two zeros of the integral \[F_n=\int_0^1\left(1-w^2\right)^n\cos(\eta w)~\mathrm dw\] as a function of $\eta$ are given in the table below.}{0.90}

Hint: The first two zeros of the integral \[F_n=\int_0^1\left(1-w^2\right)^n\cos(\eta w)~\mathrm dw\] as a function of $\eta$ are given in the table below.


The first visibility resets correspond to $\eta_1$ and $\eta_2$ for $n = 0.5$. Thus: $$\eta_2 - \eta_1 = \cfrac{2\pi}{\lambda}(d_2 - d_1)\theta_R$$$$\theta_R = \cfrac{(\eta_2 - \eta_1)\lambda}{2\pi(d_2 - d_1)} = 1.18\cdot10^{-7}$$$$R = \theta_Rh = 7.16\cdot10^8 \text{km}$$

\par\smallskip\noindent\textbf{Answer:} \textbackslash{}textbf{Answer:} $$\theta_R = \cfrac{(\eta_2 - \eta_1)\lambda}{2\pi(d_2 - d_1)} = 1.18\cdot10^{-7}$$$$R = \theta_Rh = 7.16\cdot10^8 \text{km}$$
\end{solcard}

\end{document}